\chapter{Pemetaan}\label{chapter:maps}


Kita telah menggunakan \Sage{} untuk mendefinisikan ruang vektor.
Selanjutnya, kita akan menjelajahi fungsi antar-ruang vektor.

Seperti sebelumnya, walaupun kita menganggap ruang vektor memiliki skalar
bilangan riil, di sini kita akan mendefinisikan matriks dengan entri bilangan
rasional karena entri tersebut sudah cukup untuk menggambarkan gagasannya dan
lebih mudah dibaca.
  

%========================================
\section{Fungsi linear} \label{sec:linearfunctions}
\Sage{} dapat bekerja dengan fungsi.
\begin{equation*}
  t(\colvec{a \\ b})
  =\colvec{a+2b \\ a+2b}
\end{equation*}
\begin{sagecommandline}
sage: a, b = var('a, b')
sage: t_symbolic(a, b) = [a+2*b, a+2*b]         
sage: t_symbolic       
sage: t_symbolic(1,3)       
\end{sagecommandline}

\Sage{} dapat mencari kelipatan skalar suatu fungsi atau menjumlahkan dua fungsi.
\begin{sagecommandline}
sage: t_symbolic(a, b) = [a+2*b, a+2*b]         
sage: f_symbolic = 3 * t_symbolic
sage: f_symbolic
sage: f_symbolic(4, 2)
sage: s_symbolic(a,b) = [a+b, a-b]
sage: g_symbolic = s_symbolic + t_symbolic
sage: g_symbolic
sage: g_symbolic(1, 5) 
\end{sagecommandline}
\Sage{} juga dapat mencari komposisi dua fungsi.
\begin{sagecommandline}
sage: s_symbolic * f_symbolic
sage: s_symbolic * f_symbolic (1,2)
\end{sagecommandline}

Kita memakai nama yang tidak biasa,
\inlinecode{t_symbolic}, \inlinecode{s_symbolic}, dan seterusnya,
karena objek-objek ini belum merupakan fungsi: kita belum menetapkan domain
dan kodomainnya. Objek tersebut lebih tepat disebut purwarupa fungsi.
Kita dapat memakainya untuk menghasilkan fungsi, misalnya fungsi linear
$\map{t}{\Re^2}{\Re^2}$.
\begin{sagecommandline} 
sage: a, b = var('a, b')   
sage: t_symbolic(a, b) = [a+2*b, a+2*b]         
sage: t = linear_transformation(QQ^2, QQ^2, t_symbolic)
\end{sagecommandline}
\Sage{} memakai istilah `transformasi linear' untuk semua pemetaan linear,
suatu penggunaan yang lazim, sedangkan buku ini membatasinya pada pemetaan
yang domain dan kodomainnya sama.

Untuk transformasi~$t$, kita mengharapkan
\begin{equation*}
  t(\colvec{1 \\ 3})
  =\colvec{7 \\ 7}
\end{equation*}
dan \Sage{} memberikan hasil tersebut.
\begin{sagecommandline}
sage: v = vector(QQ, [1, 3])
sage: t(v)
\end{sagecommandline}



%========================================
\section{Kiri/kanan} \label{sec:leftright}
Secara bawaan, \Sage{} merepresentasikan pemetaan linear dengan cara yang
berbeda dari buku ini.
Perbedaan tersebut paling mudah dijelaskan melalui contoh.
Tinjau kembali transformasi $\map{t}{\Re^2}{\Re^2}$.
\begin{equation*}
  t(\colvec{a \\ b})
  =\colvec{a+2b \\ a+2b}
\end{equation*}
Untuk merepresentasikannya terhadap $B,D\subset\Re^2$, kita dapat memakai
basis kanonik $B=\stdbasis_2$ dan $D=\stdbasis_2$.
Kita mencari pengaruh~$t$ pada unsur-unsur~$B$, lalu merepresentasikan
hasilnya terhadap~$D$.
\begin{equation*}
  \colvec{1 \\ 0}\mapsunder{t}\colvec{1 \\ 1}
  \text{\ menghasilkan\ }
  \rep{\colvec{1 \\ 1}}{D}=\colvec{1 \\ 1}
  \qquad
  \colvec{0 \\ 1}\mapsunder{t}\colvec{2 \\ 2}
  \text{\ menghasilkan\ }
  \rep{\colvec{2 \\ 2}}{D}=\colvec{2 \\ 2}
\end{equation*}
Dengan demikian, buku ini memperoleh matriks berikut.
\begin{equation*}
  % \hspace*{3em}
  \rep{t}{B,D}=
  \begin{mat}
    1  &2  \\
    1  &2
  \end{mat}
\end{equation*}
Namun, \Sage{} menampilkan matriks yang berbeda.
\begin{sagecommandline}
sage: t_symbolic(a, b) = [a+2*b, a+2*b]         
sage: t = linear_transformation(QQ^2, QQ^2, t_symbolic)  # lihat pembahasan 'side' di bawah
sage: t
\end{sagecommandline}
Apa yang terjadi?

Ingat bahwa tujuan representasi adalah memakai representasi matriks dari~$t$
dan representasi vektor dari~$\vec{v}$ untuk menghitung representasi vektor
dari~$t(\vec{v})$.
Buku ini menulis
\begin{equation*}
  \rep{t}{B,D}\,\rep{\vec{v}}{B}=\rep{t(\vec{v})}{D}
\end{equation*}
dengan vektor di sebelah kanan matriks.
Namun, secara bawaan \Sage{} menempatkan vektor representasi di sebelah kiri.
Perbedaannya hanya bersifat penulisan\Dash dibandingkan dengan~$T\vec{v}$
dalam buku, bentuk bawaan \Sage{} adalah
~$\trans{\vec{v}}\,\trans{T}$\Dash dan itulah sebabnya matriks \Sage{}
merupakan transpos matriks buku.

Agar \Sage{} mengikuti konvensi buku, kita akan melakukan dua hal:
mendefinisikan pemetaan dengan opsi \inlinecode{side='right'} dan menampilkan
matriksnya dengan opsi yang sama.
% This demonstrates.
\begin{sagecommandline}
sage: t_symbolic(a, b) = [a+2*b, a+2*b]         
sage: t = linear_transformation(QQ^2, QQ^2, t_symbolic, side='right')  
sage: t.matrix(side='right')
\end{sagecommandline}


  

%========================================
\section{Mendefinisikan pemetaan linear}
Kita akan menjelaskan dua cara untuk mendefinisikan pemetaan linear.

\subsection{Secara simbolis}
Kita telah memperkenalkan cara mendefinisikan pemetaan dengan rumus.
\begin{sagecommandline}
sage: a, b = var('a, b')   
sage: h_symbolic(a, b) = [a+b, a-b, b]         
sage: h_symbolic       
sage: h = linear_transformation(QQ^2, QQ^3, h_symbolic, side='right')
sage: h.matrix(side='right') 
\end{sagecommandline}
Mengevaluasi fungsi ini pada suatu anggota domain menghasilkan anggota kodomain.
\begin{sagecommandline}
sage: v = vector(QQ, [1, 3])  
sage: h(v)
\end{sagecommandline}

Jika vektor tersebut bukan anggota domain,
\begin{lstlisting}
sage: v1 = vector(QQ, [1,2,3,4])
sage: h(v1)
\end{lstlisting}
\Sage{} memberikan pesan galat pelacakan balik yang panjang. Seperti biasa,
baris terakhir paling informatif karena memuat:
\inlinecode{Type Error:(1, 2, 3, 4) fails to convert into the map's domain}.

\Sage{} dapat mencari ruang nol dan daerah hasil dengan memakai istilah
`kernel' dan `image'.
\begin{sagecommandline}
sage: h.kernel()                                       
sage: h.image()                                        
\end{sagecommandline}
\Sage{} menemukan bahwa ruang nol adalah subruang trivial dari domain,
$\nullspace{h}=\set{\zero}$.
Subruang ini berdimensi~$0$, sehingga \Sage{} melaporkan bahwa basisnya kosong.
\Sage{} juga menemukan bahwa daerah hasil $\rangespace{h}$ berdimensi~$2$
dan memberikan basisnya.

Karena ruang nol trivial, pemetaan tersebut injektif.
Karena daerah hasil berdimensi~$2$ sedangkan kodomain berdimensi~$3$,
pemetaan tersebut tidak surjektif.
Perhatikan bahwa temuan ini sesuai dengan teorema bahwa jumlah dimensi ruang
nol dan daerah hasil, yaitu $0$ dan~$2$, sama dengan dimensi domain.

Sebagai contoh lain, tinjau fungsi berikut.
\begin{sagecommandline}
sage: t_symbolic(a, b) = [a+2*b, a+2*b]
sage: t_symbolic
sage: t = linear_transformation(QQ^2, QQ^2, t_symbolic, side='right')
\end{sagecommandline}
Fungsi ini tidak injektif karena terdapat dua masukan yang dipetakan ke
keluaran yang sama.
\begin{sagecommandline}
sage: t(vector(QQ, [2, 0]))
sage: t(vector(QQ, [0, 1]))
\end{sagecommandline}
\noindent
Cara lain untuk mengetahui bahwa pemetaan tersebut tidak injektif adalah
menghitung bahwa $\nullspace{t}$ tidak berdimensi~$0$.
\begin{sagecommandline}
sage: t.kernel()
\end{sagecommandline}

Sebelum melihat daerah hasil pemetaan ini, kita telah mengetahui bahwa
dimensinya harus~$1$, sebab jumlah dimensi ruang nol dan daerah hasil sama
dengan dimensi domain.
\Sage{} mengonfirmasinya.
\begin{sagecommandline}
sage: t.image()
\end{sagecommandline}



\subsection{Melalui matriks}
Cara lain untuk mendefinisikan pemetaan linear adalah dengan matriks.
\begin{sagecommandline}
sage: M = matrix(QQ, [[1, 2], [3, 4], [5, 6]])
sage: M
sage: m = linear_transformation(M, side='right')
\end{sagecommandline}
Karena matriks tersebut memiliki entri rasional, \Sage{} mengambil domain
dan kodomainnya sebagai $\Q^2$ dan $\Q^3$.
Sekali lagi, secara bawaan \Sage{} memilih representasi dengan vektor
mengalikan dari kiri, sehingga kita menetapkan \inlinecode{side='right'}.
\begin{sagecommandline}
sage: m 
\end{sagecommandline}
Walaupun kita mendefinisikan pemetaan linear~$m$ dengan
\mbox{\inlinecode{side='right'},} matriks yang ditampilkan \Sage{} secara
bawaan menggunakan \inlinecode{side='left'}.
Untuk memperoleh representasi yang sesuai dengan buku, mintalah secara eksplisit.
\begin{sagecommandline}
sage: m.matrix(side='right')
\end{sagecommandline}

Ketika didefinisikan dengan cara ini, \Sage{} menganggap pemetaan linear
tersebut direpresentasikan oleh matriks terhadap basis standar.
\begin{sagecommandline}
sage: m = linear_transformation(M, side='right')
sage: v = vector(QQ, [7, 8])
sage: v
sage: m(v)  
\end{sagecommandline}
(Jika \inlinecode{side='right'} dihilangkan dari definisi pemetaan linear,
lalu Anda mengalikan dengan vektor beranggota dua, Anda akan menerima pesan
galat yang menyatakan bahwa vektor tersebut tidak berada dalam domain
berdimensi~$3$ dari pemetaan.)

Walaupun pemetaan berikut didefinisikan dengan dua cara, \Sage{} dapat
menggabungkannya.
\begin{sagecommandline}
sage: M = matrix(QQ, [[1, 2], [3, 4], [5, 6]])
sage: m = linear_transformation(M, side='right')
sage: n_symbolic(a, b) = [a+2*b, 3*a+4*b, 5*a+6*b]
sage: n = linear_transformation(QQ^2, QQ^3, n_symbolic, side='right')
sage: m == n  
\end{sagecommandline}
Anda juga dapat mengajukan pertanyaan yang sama pada pemetaan linear yang
dibuat dari matriks seperti yang diajukan pada pemetaan linear dari fungsi simbolis.
\begin{sagecommandline}
sage: M = matrix(QQ, [[1, 2], [3, 4], [5, 6]])
sage: m = linear_transformation(M, side='right')
sage: m.kernel() 
\end{sagecommandline}

Anda juga dapat mendefinisikan pemetaan dengan memulai dari matriks yang
merepresentasikan pemetaan terhadap basis tak standar.
Di sini kita mendefinisikan basis tak standar bagi domain dan kodomain,
$B=\sequence{\vec{\beta}_1,\vec{\beta}_2}$ dan
~$D=\sequence{\vec{\delta}_1,\vec{\delta}_2}$.
\begin{sagecommandline}
sage: M = matrix(QQ, [[1, 2], [3, 4]])
sage: beta_1 = vector(QQ, [1, -1])
sage: beta_2 = vector(QQ, [1, 1])
sage: domain_basis = [beta_1, beta_2]
sage: D = (QQ^2).subspace_with_basis(domain_basis)
sage: delta_1 = vector(QQ, [2, 0])
sage: delta_2 = vector(QQ, [0, 3])
sage: codomain_basis = [delta_1, delta_2]
sage: C = (QQ^2).subspace_with_basis(codomain_basis)
sage: m = linear_transformation(D, C, M, side='right')
sage: m.matrix(side='right')
sage: m(vector(QQ, [1, 0]))
\end{sagecommandline}
\Sage{} telah menghitung bahwa
\begin{equation*}
  \colvec{1 \\ 0}=(1/2)\colvec{1 \\ -1}+(1/2)\colvec{1 \\ 1}
  \quad\text{sehingga}\quad
  \rep{\colvec{1 \\ 0}}{B}=\colvec{1/2 \\ 1/2} 
\end{equation*}
dan menggunakannya untuk menghitung vektor citra.
\begin{equation*}
  \begin{mat}
    1 &2 \\
    3 &4
  \end{mat}
  \colvec{1/2 \\ 1/2}
  =
  \colvec{3/2 \\ 7/2}
  =
  \rep{m(\vec{v})}{D}
  \quad\text{sehingga}\quad
  m(\vec{v})=(3/2)\colvec{2 \\ 0}+(7/2)\colvec{0 \\ 3}
  =\colvec{3 \\ 21/2}
\end{equation*}





%========================================
\section{Operasi}

Tetapkan suatu ruang vektor domain~$D$ dan kodomain~$C$, lalu tinjau himpunan
semua transformasi linear di antara keduanya.
Koleksi ini memiliki beberapa operasi alami, termasuk penjumlahan dan
perkalian skalar.
Sebelumnya kita melihat bahwa \Sage{} dapat bekerja dengan operasi-operasi
tersebut serta komposisi. Kita menutup bagian ini dengan menunjukkan
pengaruhnya pada representasi.


\subsection{Penjumlahan dan perkalian skalar}
Ingat bahwa penjumlahan matriks didefinisikan sedemikian sehingga representasi
jumlah dua transformasi linear adalah jumlah matriks dari representasi
masing-masing. \Sage{} dapat menggambarkannya.
\begin{sagecommandline}
sage: M = matrix(QQ, [[1, 2], [3, 4]])
sage: m = linear_transformation(QQ^2, QQ^2, M, side='right')
sage: N = matrix(QQ, [[5, -1], [0, 7]])
sage: n = linear_transformation(QQ^2, QQ^2, N, side='right')
sage: m.matrix(side='right')
sage: n.matrix(side='right')
sage: (m+n).matrix(side='right')
\end{sagecommandline}
Tanda kurung pada baris terakhir diperlukan karena jika kita memasukkan
\inlinecode{m+n.matrix(side='right')}, \Sage{} mencoba menggabungkan pemetaan
linear~\inlinecode{m} dengan matriks \inlinecode{n.matrix(side='right')}.

Demikian pula, perkalian skalar suatu pemetaan linear tercermin dalam
perkalian skalar matriksnya.
\begin{sagecommandline}
sage: M = matrix(QQ, [[1, 2], [3, 4]])
sage: m = linear_transformation(QQ^2, QQ^2, M, side='right')
sage: (3*m).matrix(side='right')
sage: (m*3).matrix(side='right')
\end{sagecommandline}



\subsection{Komposisi}
Komposisi pemetaan linear menghasilkan perkalian matriks.
\Sage{} memakai lambang \inlinecode{*} untuk menyatakan komposisi pemetaan linear.
\begin{sagecommandline}
sage: M = matrix(QQ, [[1, 2], [3, 4]])
sage: m = linear_transformation(QQ^2, QQ^2, M, side='right')
sage: N = matrix(QQ, [[5, -1], [0, 7]])
sage: n = linear_transformation(QQ^2, QQ^2, N, side='right')
sage: M*N
\end{sagecommandline}

Seperti ditekankan dalam buku, tujuan perkalian matriks adalah
merepresentasikan komposisi pemetaan linear.
\begin{sagecommandline}
sage: (m*n).matrix(side='right')
\end{sagecommandline}




\endinput


TODO:
